% Pista ana-23j2-q1 · Anàlisi
% Procedència: PAU Catalunya 2023 · Convocatòria de juny · Sèrie 5 · Qüestió 1
\documentclass[11pt,a4paper]{article}
\usepackage{pista}
\pistainfo{Catalunya 2023}{Convocatòria de juny}{Sèrie 5}

\begin{document}

\section*{\textcolor{blauPista}{Qüestió 1}}

\noindent Considereu les funcions $f(x) = -x^{2} + x + 6$ i $g(x) = -9x + 3x^{2}$.

\subsection*{\textit{a)} \normalfont Calculeu l'àrea de la regió delimitada per les dues funcions. \hspace*{\fill}\textnormal{[1,25~punts]}}

\pista{En primer lloc, iguala $f(x) = g(x)$ i resol l'equació de segon grau: les dues solucions et marquen precisament l'interval sobre el qual integres. Per exemple, si les solucions de l'equació de segon grau fossin $x = 3$ i $x = 7$, vol dir que has d'integrar en l'interval $[3, 7]$.

\medskip
A continuació, calcula l'àrea amb el valor absolut de la integral de la diferència:
\[
\text{Àrea} = \left| \int_{x_{1}}^{x_{2}} \bigl(f(x) - g(x)\bigr)\,dx \right|
\]
Fes la resta $f(x) - g(x)$ amb cura, perquè $g$ té signes negatius. Obtindràs un polinomi de segon grau, que serà senzill d'integrar.}

\subsection*{\textit{b)} \normalfont Trobeu l'equació de la recta tangent a la funció $f(x)$ en el punt $(-2, 0)$. Representeu aquesta recta tangent i les funcions $f(x)$ i $g(x)$ en uns mateixos eixos de coordenades. \hspace*{\fill}\textnormal{[1,25~punts]}}

\pista{\textit{La tangent.} El pendent de la recta tangent que busques és $f'(-2)$, i el punt és $(-2, 0)$. Per tant:
\[
y - f(-2) = f'(-2)\,(x + 2)
\]

\medskip
\textit{El dibuix.} Les funcions $f$ i $g$ són paràboles: per dibuixar-les, n'hi ha prou amb el vèrtex i els talls amb l'eix d'abscisses. El vèrtex és el punt on s'anul·la la derivada. La recta, la pots dibuixar amb dos punts: el de tangència i el tall amb l'eix d'ordenades.}

\end{document}
